\documentclass[11pt]{article}

\usepackage[margin=1in]{geometry}
\usepackage[brazilian]{babel}
\usepackage{amsmath,amssymb,amsthm,mathtools}
\usepackage{enumitem}
\usepackage{hyperref}
\hypersetup{colorlinks=true,linkcolor=blue,citecolor=blue,urlcolor=blue}
\usepackage{orcidlink}
\usepackage{fontawesome5}

\newcommand{\ZZ}{\mathbb{Z}}
\newcommand{\NN}{\mathbb{N}}

\title{Honestidade de Escopo na Literatura Adjacente a Collatz:\\
Uma Calibração Contínua Contra Alegações de Prova Completa}

\author{Renato Augusto Tavares~{\small\href{mailto:dr.renatotavares@gmail.com}{\textcolor{black}{\faEnvelope}}}~\orcidlink{0009-0002-0196-3311}\\
\normalsize Universidade Federal de Goiás}
\date{\today}

\begin{document}

\maketitle

\begin{abstract}
Uma nota irmã deste levantamento cataloga doze alegações de prova ou
refutação completa da Conjectura de Collatz, todas inválidas. Esta nota
se volta a uma população diferente: dezesseis papers, encontrados no
mesmo levantamento bibliográfico mais amplo, que não fazem tal alegação
--- resultados sobre estruturas adjacentes a Collatz, heurísticas ou
estatísticas parciais. Lidos na íntegra e verificados computacionalmente
quando possível, esse conjunto se comporta de forma muito diferente do
conjunto de alegações de prova. O eixo que os distingue não é a taxa de
erro, mas a honestidade de escopo: sem exceção conhecida nesta amostra,
esses papers rotulam corretamente uma conjectura como conjectura e um
teorema como teorema, em contraste nítido com o overclaim sistemático
que define o catálogo irmão. Erros genuínos ocorrem --- mais
notavelmente dois num programa de sete papers de um único autor
prolífico, um num resultado
autônomo por outro lado cuidadoso --- mas se distribuem de forma
diferente: os erros do autor prolífico têm ambos a mesma forma (uma
alegação conjectural ou preliminar inconsistente com um enunciado já
estabelecido em outro lugar do mesmo texto, não um deslize aritmético)
e estão contidos a uma única seção cada; o erro do resultado autônomo
está no seu teorema titular. Esta nota é
explicitamente um documento parcial e contínuo: consolida dezesseis
itens revisados em profundidade até agora, de uma coleção levantada
muito maior, e é estruturada para crescer conforme mais itens forem
processados.
\end{abstract}

\section{Introdução}
\label{sec:intro}

Uma nota irmã, \emph{Errors in Claimed Complete Proofs (and Disproofs)
of the Collatz Conjecture}, cataloga doze alegações de solução completa
da conjectura, encontradas ao longo de um levantamento bibliográfico
mais amplo, e mostra que as doze falham --- a maioria pelo mesmo
punhado de erros recorrentes. Aquela população é incomum: todo item
nela alega explicitamente resolver a conjectura, uma alegação forte que
convida escrutínio forte.

Esta nota é sobre o restante do levantamento: papers que estudam
estruturas adjacentes a Collatz --- reformulações da dinâmica
acelerada, modelos estatísticos de comportamento de órbita,
propriedades combinatórias de sequências de paridade, sistemas de
coordenadas para o espaço de estados --- sem alegar provar ou refutar a
própria conjectura. É uma população muito maior e mais heterogênea que
o conjunto de alegações de prova, e esta nota não tenta cobri-la
inteira: nesta versão, consolida dezesseis itens que foram lidos na
íntegra e revisados em profundidade, de um conjunto consideravelmente
maior encontrado no levantamento. É escrita como um documento contínuo,
com expectativa de crescer conforme mais da coleção levantada for
processada em sessões futuras, não como um censo terminado.

A pergunta que esta nota faz sobre os dezesseis itens que cobre não é
principalmente `estão corretos' --- a maioria está, e correção em
matemática elementar ou padrão aplicada com competência não é
surpreendente. A pergunta mais informativa, e a que permite a esta nota
funcionar como ponto de calibração contra o catálogo irmão, é
honestidade de escopo: o paper alega só o que estabelece? A
Seção~\ref{sec:rc} examina o bloco coerente mais amplo desta amostra,
um programa de sete papers de um único autor aplicando um conceito
clássico de teoria ergódica por paper à mesma quantidade subjacente. A
Seção~\ref{sec:others} cobre os nove itens restantes, independentes
entre si e do programa da Seção~\ref{sec:rc}. A Seção~\ref{sec:discussion}
traça o contraste com o catálogo irmão e afirma o que a contenção de
erro generaliza e o que não generaliza. A Seção~\ref{sec:conclusion}
conclui.

\section{Metodologia}
\label{sec:methodology}

Cada item foi lido na íntegra, e suas alegações concretas e
verificáveis --- identidades, fórmulas fechadas, teoremas com enunciado
explícito --- foram verificadas computacionalmente de forma
independente sempre que admitiam um teste numérico. Onde um paper
distingue resultados provados de conjecturais, essa rotulagem foi
conferida contra o que a própria demonstração do paper de fato
estabelece, não aceita pelo valor de face. Nos poucos casos em que o
julgamento envolvido era sutil, a análise foi cruzada com um revisor
independente baseado em modelo de linguagem. Como no catálogo irmão,
toda alegação específica feita abaixo sobre um paper nomeado remonta a
esse processo de revisão, não a uma paráfrase dele; o detalhe técnico
completo de cada item foi registrado durante a revisão em notas de
trabalho internas e é resumido, não reproduzido, aqui.

\section{O programa Ruiz Castillo}
\label{sec:rc}

Sete papers de Juan Carlos Ruiz Castillo (ScD., Universidad de San
Carlos de Guatemala) aparecem nesta amostra, cada um aplicando um
conceito clássico de teoria ergódica ou de grandes desvios --- pressão,
uma `geometria' residual, cotas dissipativas, um teorema central do
limite, um operador de transferência, medidas de Gibbs, um princípio
variacional, uma função de tasa de grandes desvios --- à mesma
quantidade elementar: o drift residual $L_k(n) = k\log_2 3 - A_k(n)$,
onde $A_k(n)$ é a soma das primeiras $k$ valuações 2-ádicas ao longo do
mapa acelerado de Collatz $U(n) = (3n+1)/2^{v_2(3n+1)}$ --- o drift
padrão da equação de ciclo, já usado em outras partes deste programa de
pesquisa. Nenhum dos sete alega provar a Conjectura de Collatz; cada um
diz isso explicitamente. A apresentação é fortemente repetitiva ao
longo da série (quadros decorativos recorrentes redesenhando a mesma
cadeia conceitual, taxas de autocitação de 75--100\% revelando da ordem
de 19--20 títulos quase idênticos do mesmo autor), mas o rigor formal
subjacente é notavelmente estável: cinco dos sete papers não contêm
nenhum erro que encontramos, e os dois erros genuínos encontrados estão
ambos contidos a uma única seção e não afetam as conclusões
posteriores dos próprios papers.

\textbf{Geometría Residual} \cite{RuizCastillo2026a} aplica o
formalismo termodinâmico padrão (pressão como supremo sobre medidas
invariantes, uma segunda derivada tipo métrica de Fisher) ao drift
residual, corretamente e sem alegar provar Collatz. O mesmo fato
trivial --- convexidade implica segunda derivada não-negativa --- é
apresentado como cinco resultados nomeados separadamente ao longo do
paper, emoldurado por dez quadros decorativos redesenhando a mesma
cadeia conceitual, sem nenhum conteúdo numérico em lugar nenhum (ambas
as figuras são explicitamente rotuladas `conceptual') e 100\% de
autocitação (13 de 13 referências). Calculamos explicitamente, em forma
fechada, a quantidade que o paper deixa abstrata: sob o modelo i.i.d.
padrão já usado neste programa de pesquisa, $\Lambda(t) = t(1-\log_2 3)
- \log(2-e^t)$, cuja segunda derivada diverge quando $t \to \log 2^-$
--- uma característica genuína, ainda que incidental, que o paper nunca
calcula porque nunca deriva $\Lambda(t)$ explicitamente.

\textbf{Dissipative Bounds} \cite{RuizCastillo2026b}, o segundo paper
deste autor nesta amostra, repete o padrão: identidades elementares
corretas (um `desenrolar' afim exato do mapa acelerado, uma cota de
dissipação universal), sem conteúdo numérico, e autocitação de 75\%
(12 de 16 referências), com a própria lista de autocitação revelando
mais cerca de uma dúzia de títulos quase idênticos do mesmo autor sobre
a mesma `decomposição residual'.

\textbf{Teorema Central del Límite Residual} \cite{RuizCastillo2026c}
rotula sua alegação central com honestidade: o título diz `Teorema',
mas o resultado é enunciado no corpo como `Conjetura 4.2', com cinco
hipóteses explícitas não estabelecidas (existência de uma medida de
Gibbs residual, ergodicidade, um espaço funcional adequado, uma brecha
espectral, variância residual positiva), e a própria conclusão do paper
afirma claramente que esse arcabouço não prova a Conjectura de Collatz.
Testamos a previsão observável da conjectura --- normalidade
assintótica das flutuações do drift residual --- diretamente em órbitas
reais de Collatz (não no modelo i.i.d. abstrato): a variância se
estabiliza perto de 2, a assimetria tende a 0 e a curtose tende a 3
conforme a janela cresce, um resultado que o próprio paper nunca testa
numericamente.

\textbf{Operador de Transferencia Residual} \cite{RuizCastillo2026d}
contém o primeiro erro genuíno deste programa. Sua Proposición 5.3
alega $\lim_{t\to\infty} L_t(1) = 0$; sua própria demonstração deriva a
forma fechada $L_t(1) = e^{(\log_2 3 - 1)t}/(1-e^{-t})$ e observa
corretamente que isso `cresce exponencialmente quando $t\to\infty$',
imediatamente seguido pelo símbolo de fim de demonstração, sem que a
alegação `$=0$' seja corrigida. Confirmamos numericamente que $L_t(1)$
cresce de 2,84 em $t=1$ para cerca de $6,4\times 10^{50}$ em $t=200$,
estritamente crescente ao longo de toda a faixa --- o oposto exato do
limite alegado. Isso é uma inconsistência entre enunciado e
demonstração, não um erro de cálculo: a própria demonstração deriva o
assintótico correto, e o texto ao redor usa esse assintótico correto
para motivar a normalização da seção seguinte. O erro está confinado a
um cálculo preliminar da Seção 5; nenhum resultado posterior
explicitamente conjectural (Seções 6--8, sobre o operador de
transferência restrito) depende dele.

\textbf{Medidas de Gibbs Residuales} \cite{RuizCastillo2026e} e
\textbf{Principio Variacional Residual} \cite{RuizCastillo2026f} voltam
ao padrão sem erro: reafirmações corretas da mesma identidade central,
aparato clássico de análise convexa (supremos de funções afins,
dualidade de Legendre--Fenchel) corretamente aplicado, e todo resultado
aberto honestamente rotulado `Conjetura', nunca `Teorema' ou
`Proposición'.

\textbf{Grandes Desviaciones Residuales} \cite{RuizCastillo2026g}
contém o segundo erro genuíno deste programa, estruturalmente
semelhante em tipo ao primeiro, mas localizado de forma diferente: sua
Definición 3.1 define uma função de tasa $I_{\mathrm{RC}}(x)$ via o
evento de cauda unilateral $\{L_k/k \ge x\}$, e sua Proposición 3.4
(Seção 3, já provada) mostra corretamente que essa função de tasa
unilateral é não-decrescente. Sua Seção 7, porém, inteiramente
conjectural --- Figura 1, Conjetura 7.3 e Conjetura 7.5, mutuamente
consistentes entre si --- descreve uma função de tasa de Cramér
clássica e bilateral, positiva nos dois lados de seu zero único,
contradizendo a monotonicidade já estabelecida na Seção 3. Confirmamos
a contradição por três rotas independentes (a fórmula padrão de
grandes desvios unilateral, simulação Monte Carlo e um cálculo exato
Binomial Negativo para $k$ pequeno), todas concordando que a tasa
unilateral correta é identicamente zero abaixo do drift típico, não só
num ponto. Como no caso anterior, o erro é uma inconsistência de uma
seção conjectural com uma proposição já provada mais cedo no mesmo
texto, não um engano aritmético, e está confinado à Seção 7; a
conclusão do paper (Seção 8) não depende da forma específica de
$I_{\mathrm{RC}}$.

\section{Nove papers independentes}
\label{sec:others}

Os nove itens restantes desta amostra não compartilham autoria nem,
além de estruturas adjacentes a Collatz, muito assunto entre si.

\textbf{Adnan \& Dar} \cite{AdnanDar2026} estendem a regra $3n+1$ a
decimais em $(0,1)$ definindo `paridade' via o último dígito
significativo de um número, concluindo divergência universal. Todos os
exemplos trabalhados no paper conferem aritmeticamente, mas a
construção é conceitualmente mal-posta: `último dígito significativo'
não está definido para nenhum decimal com expansão infinita (uma
exceção de medida zero na leitura ingênua, mas que exclui $1/3$, $1/7$
e todo irracional), e mesmo restrita a decimais finitos a divergência é
um artefato trivial de um termo $+1$ não normalizado, não um análogo
estrutural do mecanismo real de Collatz (onde $3n+1$ é sempre par para
$n$ ímpar por construção). Não é uma alegação de prova sobre a
conjectura original e não pertence ao catálogo irmão.

\textbf{Gilbert} \cite{Gilbert2026} constrói uma conjugação explícita
entre o grafo podado de Collatz e os inteiros não-nulos, codificando a
classe residual descartada mod 3 via sinal, negando explicitamente
qualquer prova da conjectura. Todo teorema que checamos --- a própria
conjugação, o pruning de múltiplos de 3, o mapa acelerado, a fórmula de
contagem de pais e uma conexão com a sequência OEIS A254046 --- vale
exatamente. Nenhum erro encontrado.

\textbf{Seymour (Steiner Sentence Length)} \cite{Seymour2026a} deriva
uma matriz de probabilidades de transição de um passo entre resíduos
mod 8 (Teorema 2.1), que confirmamos correta, e a usa para derivar uma
distribuição de forma fechada alegada para o comprimento de `sentenças
Steiner', $3^{k-1}/4^k$ (Teorema 5.1), rejeitando a alternativa ingênua
$(1/2)^k$ via um teste $\chi^2$ reportado. Nossa reimplementação
independente das próprias definições do paper dá o resultado oposto:
a simulação bate com $(1/2)^k$, e um contraexemplo aritmético exato
($n=68567$) mostra uma sentença de comprimento 1 cujo resíduo de
entrada é 7 --- um caso que a demonstração do Teorema 5.1 nunca
contempla (só conta sentenças cuja \emph{primeira} palavra já fecha,
deixando de fora silenciosamente palavras com entrada 3 ou 7 que podem
fechar já no primeiro passo). Reportamos isso com cautela real: o paper
alega uma formalização completa em Lean 4/Mathlib sem \texttt{sorry}, e
não temos acesso a esse arquivo de formalização; é possível que a prova
Lean verifique corretamente uma definição que diverge sutilmente do
enunciado em prosa, caso em que o problema estaria na tradução
prosa-para-formalização, não na própria prova formal. Com base
estritamente nas definições como escritas no paper, porém, tanto um
contraexemplo exato quanto uma simulação em larga escala discordam do
resultado central.

\textbf{Seymour (Regular Expression Language)} \cite{Seymour2026b}, um
`working paper' explicitamente rotulado como tal pelo mesmo autor,
separa teoremas provados de conjecturas ao longo de todo o texto. Dois
resultados centrais (Proposición 3.1 sobre a aritmética do circuito
Steiner, Teorema 3.6 sobre comprimento de corridas) conferiram
exatamente em teste de larga escala. Um erro pequeno e contido:
Corolario 2.2 alega que um resíduo de saída mod 24 depende da paridade
de uma variável auxiliar $t$; o teste mostra que o valor alegado só
vale quando $t$ é par, e o resíduo real é a mesma constante
independentemente da paridade de $t$, contradizendo a alegação do
próprio paper de que o mapa $t \mapsto (3t+1)/2$ preserva paridade
(falso em geral). Esse corolário não consta na própria lista do paper
de resultados provados, e o erro não afeta a conjectura central do
working paper, que decorre quase imediatamente do já verificado Teorema
2.1 do paper companheiro.

\textbf{Anthony} \cite{Anthony2026} constrói um aparato formal elaborado
(equação P de Riemann, monodromia, identidades hipergeométricas) em
torno de uma reformulação de dois campos da dinâmica de Collatz via a
coordenada de Möbius $\Phi(n) = 1/(n+1)$. Toda alegação verificável ---
a reformulação de Möbius, um teorema excluindo a persistência
indefinida do regime $v_2(3n+1)=1$, e várias identidades padrão da
função digamma --- conferiu exatamente. O paper é incomumente cuidadoso
ao longo de todo o texto em separar teoremas provados de analogias
estruturais (explicitamente: `this is an analogy of form, not of
kind') e de resultados condicionais cujas condições não estão
estabelecidas, e corrige, dentro do próprio texto, uma alegação
anterior equivocada de uma versão prévia da mesma linha de pesquisa.

\textbf{Hikawa} \cite{Hikawa2026} estuda estruturas combinatórias
finitas de vetores de paridade, independentemente de a Conjectura de
Collatz ser verdadeira ou falsa --- uma ressalva que o paper repete em
quase toda seção. Ambos os teoremas centrais que testamos (uma bijeção
explícita entre classes de vetores, e uma propriedade de rigidez da
valuação 2-ádica de um termo de correção) valeram exatamente por força
bruta, incluindo o próprio caso de exceção declarado pelo paper.

\textbf{Olgac} \cite{Olgac2026} é uma nota de três páginas conectando
duas alegações anteriores não catalogadas do mesmo autor, uma sobre
Collatz e outra alegando uma prova completa da Conjectura de Goldbach
(fora do escopo deste levantamento). Seu único conteúdo específico de
Collatz é uma tautologia (ramos raiz-disjuntos da árvore reversa,
consequência imediata da fatoração única) em vez de um resultado
genuíno; o restante da nota não é enunciado com precisão matemática
suficiente para avaliar como correto ou incorreto.

\textbf{Williams} \cite{Williams2026} organiza os inteiros positivos por
sua fatoração 3-suave e estuda a dinâmica de Collatz como um fluxo
diagonal determinístico dentro dessa organização. Todo teorema que
testamos --- uma identidade da dinâmica diagonal, um teorema sobre
linhas livres de primos, uma partição bijetora, uma contagem
assintótica de posições admissíveis --- valeu exatamente. Dois erros
pequenos e autocontidos em citações OEIS foram encontrados (um número
de sequência que não bate com a propriedade citada, uma descrição
com erro de um elemento), nenhum afetando qualquer prova do paper. O
paper separa teoremas provados de observações computacionais e
problemas em aberto explícitos ao longo de todo o texto, e declara seu
uso de IA generativa (para rascunho e formatação, com verificação
humana integral declarada).

\textbf{Fu, Liu \& Wang} \cite{FuLiuWang2026} propõem um mecanismo de
processo de Poisson explicando uma estatística tipo Gama previamente
observada para comprimentos de órbita de Collatz. Toda alegação
verificável --- a esperança e a variância da valuação 2-ádica, uma
aproximação de forma fechada ligando comprimento de órbita a um valor
inicial, e a álgebra de uma condição de fechamento de ciclo --- valeu
exatamente. O paper é explícito, corretamente e em contraste com a
tentativa de prova falha revisada no catálogo irmão sob o mesmo estilo
geral de argumento, que sua obstrução assintótica de fechamento de
ciclo não é uma prova de que nenhum ciclo não-trivial existe.

\section{Discussão}
\label{sec:discussion}

O contraste que se mantém em toda esta amostra, sem exceção conhecida,
é a honestidade de escopo: nenhum destes dezesseis papers veste uma
conjectura de roupa de prova da própria Conjectura de Collatz, mesmo
onde a matemática subjacente é elementar e fortemente repetida (o
programa Ruiz Castillo) e mesmo onde um resultado central é
provavelmente errado (o paper de sentenças Steiner do Seymour). Esse é
o eixo em que esta população difere dos doze casos do catálogo irmão,
onde todo item sem exceção faz a alegação forte e todo item sem exceção
falha em sustentá-la. Taxa de erro não é o traço que distingue aqui ---
ambas as populações contêm erros reais --- overclaim é.

Contenção de erro é uma observação mais estreita, mas verdadeira, só
que vale sobre o programa Ruiz Castillo especificamente, não como
propriedade desta população como um todo. Os dois erros genuínos
daquele programa têm a mesma forma (uma alegação conjectural ou
preliminar inconsistente com um enunciado já estabelecido em outro
lugar do mesmo texto, não um deslize aritmético) e ambos estão
confinados a uma única seção, deixando intocado o material posterior e
explicitamente conjectural do paper. O paper de sentenças Steiner do
Seymour é o contraponto nítido: seu erro está no Teorema 5.1 titular, o
resultado central alegado do paper, não num lema periférico. O fato
interessante não é que erros nesta literatura sejam contidos --- em
geral não são --- mas que o programa de autor único mais prolífico
desta amostra erra de forma consistentemente contida, enquanto o
resultado autônomo mais afiado da amostra erra em seu centro. Essa
diferença é um fato sobre como autores específicos erram, não uma
generalização sobre a população.

Notamos, por fim, que a cobertura deste capítulo é parcial por
construção: dezesseis itens revisados em profundidade aqui, extraídos
de uma coleção levantada consideravelmente maior que inclui tanto mais
papers adjacentes a Collatz ainda não processados quanto um corpo
separado de literatura de referência estabelecida (o teorema de
quase-todas-as-órbitas de Tao, o modelo estocástico de Kontorovich--Lagarias
e semelhantes) que é citado ao longo deste programa de pesquisa mas
nunca foi alvo deste tipo de revisão. Espera-se que revisões futuras
desta nota acrescentem mais itens conforme forem processados; a
estrutura acima --- uma seção por autor ou programa coerente,
parágrafos individuais por item independente --- é pensada para
acomodar esse crescimento sem exigir reorganização.

\section{Conclusão}
\label{sec:conclusion}

Dezesseis papers adjacentes a Collatz, nenhum alegando provar ou
refutar a conjectura, revisados na íntegra: nove não contêm nenhum erro
que encontramos; quatro contêm um erro pequeno e contido (dois no
programa Ruiz Castillo, um num working paper do Seymour, um par de
deslizes só de citação no Williams); um --- o paper de sentenças
Steiner do Seymour --- tem um erro provável no seu teorema central e
titular, não num detalhe periférico; e dois não fazem nenhuma alegação
avaliável diretamente como correta ou incorreta (a construção de Adnan
\& Dar é conceitualmente mal definida; o único conteúdo específico de
Collatz do Olgac é tautológico). Nos dezesseis, a distinção
teorema-vs-conjectura é mantida com um rigor que o catálogo irmão de
alegações de prova deste levantamento nunca mostra. Esse contraste ---
honestidade de escopo, não taxa de erro --- é o ponto de calibração
desta nota, e espera-se que ele se aprofunde, não mude de natureza,
conforme mais da coleção levantada for acrescentada em revisões
futuras.

\section*{Agradecimentos}

Partes da verificação descrita nesta nota (reconstrução de casos,
checagens computacionais e cruzamento de julgamentos matemáticos mais
sutis) foram realizadas com o auxílio de Claude (modelos das famílias
Claude 4 e 5, Anthropic), usado como assistente de pesquisa sob a
direção e revisão do autor.

\begin{thebibliography}{99}

\bibitem{RuizCastillo2026a} J.\ C.\ Ruiz Castillo, \emph{Geometría
  Residual de Ruiz Castillo para la dinámica acelerada de la Conjetura
  de Collatz}, Zenodo, DOI 10.5281/zenodo.21271452, 2026.

\bibitem{RuizCastillo2026b} J.\ C.\ Ruiz Castillo, \emph{Dissipative
  Bounds and Ruiz Castillo Residual Decomposition in the Accelerated
  Dynamics of the Collatz Conjecture}, ResearchGate preprint.

\bibitem{RuizCastillo2026c} J.\ C.\ Ruiz Castillo, \emph{Teorema Central
  del Límite Residual de Ruiz Castillo para la dinámica acelerada de la
  Conjetura de Collatz}, ResearchGate preprint.

\bibitem{RuizCastillo2026d} J.\ C.\ Ruiz Castillo, \emph{Operador de
  Transferencia Residual de Ruiz Castillo y estructura espectral en la
  dinámica acelerada de la Conjetura de Collatz}, ResearchGate preprint.

\bibitem{RuizCastillo2026e} J.\ C.\ Ruiz Castillo, \emph{Medidas de
  Gibbs Residuales de Ruiz Castillo y estados de equilibrio en la
  dinámica acelerada de la Conjetura de Collatz}, ResearchGate preprint.

\bibitem{RuizCastillo2026f} J.\ C.\ Ruiz Castillo, \emph{Principio
  Variacional Residual de Ruiz Castillo y formalismo termodinámico para
  la dinámica acelerada de la Conjetura de Collatz}, ResearchGate
  preprint, 2026.

\bibitem{RuizCastillo2026g} J.\ C.\ Ruiz Castillo, \emph{Grandes
  Desviaciones Residuales de Ruiz Castillo en la dinámica acelerada de
  la Conjetura de Collatz}, Zenodo, DOI 10.5281/zenodo.20767811, 2026.

\bibitem{AdnanDar2026} J.\ Adnan e S.\ A.\ Dar, \emph{Behavior of a
  Decimal-Parity-Based $3n+1$ Mapping}, Global Journal of Pure and
  Applied Mathematics, 2026.

\bibitem{Gilbert2026} J.\ S.\ Gilbert, \emph{A Collatz-Equivalent Map on
  the Nonzero Integers}, Preprints.org, DOI
  10.20944/preprints202607.0575.v1, 2026.

\bibitem{Seymour2026a} J.\ Seymour, \emph{First-Principles Derivation of
  the Steiner Sentence Length Distribution}, wildducktheories.com
  preprint, 2026.

\bibitem{Seymour2026b} J.\ Seymour, \emph{A Regular Expression Language
  for the Collatz Graph}, wildducktheories.com working paper, 2026.

\bibitem{Anthony2026} M.\ M.\ Anthony, \emph{A Two-Field Propagation
  Model for the Collatz Map: A Möbius Reformulation, Monodromy
  Structure, and a 2-adic Non-Expansion Theorem}, ResearchGate
  preprint, 2026.

\bibitem{Hikawa2026} K.\ Hikawa, \emph{Finite-Dimensional Combinatorial
  and Arithmetic Structures of Parity Vectors for the Accelerated
  Collatz Map}, ResearchGate preprint.

\bibitem{Olgac2026} E.\ Olgac, \emph{Technical Note: Structural Dualism
  in Integer Architectures}, ResearchGate preprint.

\bibitem{Williams2026} J.\ Williams, \emph{A Coordinate System for
  Collatz Dynamics}, arXiv:2607.01718, 2026.

\bibitem{FuLiuWang2026} W.\ Fu, X.\ Liu, e Y.\ Wang, \emph{Emergence of
  Gamma-Type Upward-Phase Statistics in the Collatz Map: An Effective
  Poisson Process Mechanism}, arXiv:2606.26811, 2026.

\end{thebibliography}

\end{document}
